%% Nature-Communications-style preprint template.
%% Compile: latexmk -pdf main.tex (uses bibtex via latexmkrc).
%% Engine: pdflatex by default. If authors or figures require Unicode, switch
%% to xelatex by setting `$pdf_mode = 5;` in latexmkrc and compiling with
%% `latexmk -xelatex`.
\documentclass[11pt,a4paper]{article}

\usepackage[T1]{fontenc}
\usepackage[utf8]{inputenc}
\usepackage{lmodern}
\usepackage[margin=2.1cm]{geometry}
\usepackage{graphicx}
\usepackage{amsmath,amssymb}
\usepackage{booktabs}
\usepackage{array}
\usepackage{caption}
\usepackage{subcaption}
\usepackage{microtype}
\usepackage{authblk}
\usepackage{siunitx}
\usepackage{xcolor}
\usepackage[hidelinks,colorlinks=true,citecolor=blue,linkcolor=blue,urlcolor=blue]{hyperref}
\usepackage{setspace}
%% Reference style: natbib numbered unsrtnat (Nature portfolio, BCJ, Leukemia,
%% Cell Reports Medicine, npj Precision Oncology). For author-year journals
%% (Genome Medicine, BMC), switch to
%%   \usepackage[round,authoryear,sort]{natbib}
%%   \bibliographystyle{plainnat}
%% For Vancouver style (Briefings in Bioinformatics): \bibliographystyle{vancouver}.
\usepackage[numbers,super,sort&compress]{natbib}
\usepackage{titlesec}
\usepackage{lineno}

\titleformat{\section}{\fontsize{12}{14}\bfseries\sffamily}{\thesection}{0.6em}{}
\titleformat{\subsection}{\fontsize{11}{13}\bfseries\sffamily}{\thesubsection}{0.6em}{}
\titleformat{\subsubsection}{\fontsize{10}{12}\itshape\sffamily}{}{0em}{}
\setlength{\parskip}{0.25\baselineskip}
\setlength{\parindent}{0pt}
\captionsetup{font=small,labelfont=bf,labelsep=period}
\renewcommand{\thefigure}{\arabic{figure}}
\renewcommand{\familydefault}{\rmdefault}

\title{\fontsize{15}{18}\bfseries\sffamily
Placeholder title: one concrete finding in at most 15 words}

\author[1,*]{Anonymous}
\affil[1]{Placeholder affiliation.}
\affil[*]{Correspondence: \texttt{author@example.org}.}
\date{\today}

\begin{document}

\maketitle
\thispagestyle{empty}
\vspace{-2em}

\begin{abstract}
\noindent
\textbf{Background.}\, One sentence on field state and the gap.
\textbf{Methods.}\, Pipeline summary, cohort sizes, validation strategy.
\textbf{Results.}\, Headline quantitative finding with effect size and p-value; second finding; external validation; incremental-value benchmark; honest negative if any.
\textbf{Conclusion.}\, What the paper does and does not claim; release statement.
\end{abstract}

\medskip
\noindent\textbf{Keywords:} placeholder, placeholder, placeholder.

\section{Introduction}\label{sec:intro}
Field state and clinical / biological motivation\citep{placeholderCiteA}.

Method primitive in one-sentence intuitions and foundational references\citep{placeholderCiteB}.

The gap and the concrete contribution, numbered.

\section{Results}\label{sec:results}

\subsection{Study design and cohort composition}
Describe cohorts; include age distributions, sample sizes, event counts, follow-up.

\subsection{Core method result}
\subsection{Stability and null-model results}
\subsection{Primary outcome result}
\subsection{External validation}
\subsection{Incremental value over baseline}
\subsection{Proportional-hazards diagnostics and calibration}
\subsection{Molecular or mechanistic correlates}

\section{Discussion}\label{sec:discussion}
Central finding, three caveats, three positives, limitations, outlook.

\section{Methods}\label{sec:methods}

\subsection{Data sources and harmonisation}
\subsection{Core method}
\subsection{Stability and null}
\subsection{Statistical tests}
\subsection{Reproducibility}

%% -------------------------------------------------------------------------
%% Reporting-guideline compliance
%% -------------------------------------------------------------------------
%% Clinical-prediction manuscripts: attach a TRIPOD 2015 checklist as
%% Supplementary Table ST1 (Collins et al., Ann Intern Med 2015).
%% Observational studies: STROBE checklist.
%% RCTs: CONSORT 2010.
%% Diagnostic-accuracy studies: STARD 2015.
%% Animal studies: ARRIVE 2.0.
%% Name the guideline and attach a completed checklist; reviewers check.
\subsection*{Reporting-guideline compliance}
This study follows the \textbf{<TRIPOD / STROBE / CONSORT / STARD / ARRIVE>} reporting guideline. A completed checklist is provided as Supplementary Table ST1.

\subsection*{Preregistration}
The primary outcome, primary hypothesis and pre-specified secondary outcomes were recorded in \texttt{docs/prereg.md} at commit \texttt{<hash>} prior to any analysis touching the outcomes. Exploratory analyses are labelled as such.

\subsection*{Data availability}
Public datasets with URLs, accession numbers, and a persistent DOI (Zenodo / FigShare) for processed data where licensing allows.

\subsection*{Code availability}
Source code, LaTeX sources and compiled PDF at \url{https://github.com/<user>/<slug>} under the MIT licence. The exact commit tagged \texttt{v<MAJOR>.<MINOR>.<PATCH>} reproduces the figures and tables in this manuscript. Persistent code DOI via Zenodo: \texttt{10.5281/zenodo.<id>} (created via GitHub--Zenodo integration).

\subsection*{Ethics and data-use statement}
Both datasets used here are fully de-identified publicly released resources. No new human-subjects data were collected. The work falls under the scope of public-data reanalysis and did not require additional institutional review board approval under applicable guidelines.

\subsection*{Acknowledgements}
\subsection*{Author contributions}
Following the CRediT taxonomy: \textbf{<Name>}: conceptualisation, methodology, software, formal analysis, writing -- original draft. \textbf{<Name>}: supervision, writing -- review and editing.

\subsection*{Competing interests}
The authors declare no competing interests.

\bibliographystyle{unsrtnat}
\bibliography{references}

\end{document}
